%% Lab Report for EEET2493_labreport_template.tex
%% V1.0
%% 2019/01/16
%% This is the template for a Lab report following an IEEE paper. Modified by Francisco Tovar after Michael Sheel original document.


%% This is a skeleton file demonstrating the use of IEEEtran.cls
%% (requires IEEEtran.cls version 1.8b or later) with an IEEE
%% journal paper.
%%
%% Support sites:
%% http://www.michaelshell.org/tex/ieeetran/
%% http://www.ctan.org/pkg/ieeetran
%% and
%% http://www.ieee.org/

%%*************************************************************************
%% Legal Notice:
%% This code is offered as-is without any warranty either expressed or
%% implied; without even the implied warranty of MERCHANTABILITY or
%% FITNESS FOR A PARTICULAR PURPOSE! 
%% User assumes all risk.
%% In no event shall the IEEE or any contributor to this code be liable for
%% any damages or losses, including, but not limited to, incidental,
%% consequential, or any other damages, resulting from the use or misuse
%% of any information contained here.
%%
%% All comments are the opinions of their respective authors and are not
%% necessarily endorsed by the IEEE.
%%
%% This work is distributed under the LaTeX Project Public License (LPPL)
%% ( http://www.latex-project.org/ ) version 1.3, and may be freely used,
%% distributed and modified. A copy of the LPPL, version 1.3, is included
%% in the base LaTeX documentation of all distributions of LaTeX released
%% 2003/12/01 or later.
%% Retain all contribution notices and credits.
%% ** Modified files should be clearly indicated as such, including  **
%% ** renaming them and changing author support contact information. **
%%*************************************************************************

\documentclass[journal]{IEEEtran}

% *** CITATION PACKAGES ***
\usepackage[style=ieee]{biblatex} 
\bibliography{example.bib}    

% *** MATH PACKAGES ***
\usepackage{amsmath}

% *** PDF, URL AND HYPERLINK PACKAGES ***
\usepackage{url}
% correct bad hyphenation here
\usepackage{graphicx}  %needed to include png, eps figures
\usepackage{float}  % used to fix location of images i.e.\begin{figure}[H]

\begin{document}

% paper title
\title{Title}

% author names 
\author{Student name 1}% <-this % stops a space
        

% make the title area
\maketitle

% As a general rule, do not put math, special symbols or citations
% in the abstract or keywords.

\section{Introduction }
Include a title and one or two paragraphs describing what you plan to do. Tell what interviews, site visits, or other activity you plan. Be specific if you can. Include  one good reference you plan to used. This is an example of how to include a citation \cite{williams1989style}.  \\.

\section {Background}
Give a background on your topic. Include references. 

\section{Results}
Results of your interviews or observations. Use information and/or quotes from your interview or observations.

\section{Discussion}
Your comments or evaluation of interview or observations

\section{Summary}
Summary of your paper.

\appendix
The Appendix should contain the name, position, and company (or other relevant information) for the person(s) you interviewed or the places you visited. For interviews, include your list of questions and indicate if the interview was in person, by phone, or by email. (In-person interviews are best, but may not be available for some topics.) Include the person's answers. (A summary is ok.) If you identify the person fully and quote extensively from the interview in the body of your paper you do not have to include the appendix. The Appendix does not count toward the 4000 word requirement.


\section{Summary}
Summarize your paper.




\printbibliography

% that's all folks
\end{document}


